\section{Results and Discussion}
\label{sec:results}

\subsection{Imperceptibility}
Table~\ref{tab:imperc} reports embedding distortion. Over the corpus, Variant A reaches 45.479~dB / SSIM 0.98847 and Variant C reaches 45.492~dB / 0.98841, with a minimum SSIM of 0.97831 over the whole $B=8$ set. On the author-supplied photograph the figures are 45.648~dB / 0.99356 (A) and 45.661~dB / 0.99362 (C). The two variants are practically indistinguishable in distortion, as expected from LSB shifting (about 0.5 LSB on average against about 1.5 for replacement, with tag-carrying pixels never shifted).

\subsection{Localization}
Table~\ref{tab:results} gives localization and recovery. Block recall is 1.0000 for splicing, crop-and-refill and noise, and 0.99998 for inpainting. The shortfall is 17 MSB-preserved-miss blocks in total, the predicted residual of binding the tag to the six MSB planes: a smooth inpainter can reproduce the six MSB planes exactly, and such a block is accepted by construction. Block precision is 1.0000, but this is structurally forced by the ground-truth rule, under which a block counts as tampered if the region covers any of its pixels, and is therefore not evidence of detector quality.

Pixel precision averages 0.9599 and rises from 0.9248 at $\alpha=0.10$ to 0.9797 at $\alpha=0.70$. The shortfall from 1.0 is the $8\times8$ block-quantization penalty, which weighs less as the region grows relative to its boundary; it is not a detector error and cannot shrink without a finer grid. On the photograph, block precision and block recall are both 1.0000 for every class and ratio, with pixel precision 0.9525 on average, from 0.9239 at $\alpha=0.10$ to 0.9752 at $\alpha=0.70$, for the same reason. Neighborhood refinement flagged only 2 blocks out of 10,092,544 verified.

The null condition comprises 320 rows and 1,802,240 block verifications of unmodified watermarked images, with zero false positives and a rule-of-three 95\% upper bound of $1.66\times10^{-6}$. This is a determinism and implementation-consistency check across images, variants and keys, not a measurement of the $2^{-32}$ cryptographic false-accept probability.

\subsection{Recovery}
Across the seven ratios, $\rho$ takes the values 0.9590, 0.8584, 0.7479, 0.6371, 0.5335, 0.4437 and 0.3579 (Fig.~\ref{fig:rho}). It is monotone and near-linear, and sits 3 to 6 percentage points above its structural bound $1-\alpha$ at every ratio. The mean over the main grid is 0.6482. In-region PSNR is 28.49~dB for Variant A and 31.13~dB for Variant C, essentially flat across ratios, whereas whole-image quality falls with coverage. Mapping coverage, not descriptor fidelity, is therefore the dominant constraint on recovery quality: where a partner survives, the block is rebuilt at roughly constant fidelity, and the fraction of blocks with a surviving partner is what decays.

The photograph is a deliberately harder case: its mean $\rho$ is 0.6409 and its in-region PSNR is 23.28~dB (A) and 26.56~dB (C), about 5~dB below the corpus means. It is dominated by high-frequency foliage, which a 12-coefficient DCT descriptor cannot carry, so it is an honest worst-case illustration, not a typical one. Its $\rho$ follows the corpus trajectory, falling from 0.9775 at $\alpha=0.10$ to 0.3154 at $\alpha=0.70$, so the coverage argument holds on a texture-dominated image; only per-block fidelity differs, reflecting descriptor bandwidth, not the mapping.

Figure~\ref{fig:qual} shows a splice at $\alpha=0.30$ with Variant A. Verification flags 1,292 blocks, of which 321 are unrecoverable, giving $\rho=0.7515$. Pixel precision is 0.9522 and in-region PSNR is 23.27~dB. The cyan blocks in the detection panel are marked, not fabricated.

\subsection{Comparison and Ablation}
Table~\ref{tab:compare} sets the measured figures beside published ones, which are cited and not re-implemented. Reference-sharing and fountain-coded schemes \cite{zhang2011reference,korus2013efficient} reach far higher whole-image reconstruction at high damage, for example 37~dB at 50\% damage for the latter. One-to-one allocation gives that up: a flagged block whose single partner is also flagged is lost. Its measured 17.39~dB at 50\% tamper (unmarked convention) is correspondingly lower, and 12.40~dB for the honest marked output. Published recovered qualities of 27.64~dB worst case for AuSR1 \cite{ausr1} and above 30~dB at 50\% tamper for Wu et al.\ \cite{wu2026multifeature} are likewise higher; these come from differing protocols and are not strictly commensurable, while the watermarked fidelity of 45.48~dB is comparable with the 45.57~dB reported for AuSR1. In exchange, it offers determinism, an auditable keyed construction, and disclosure of the unrecoverable fraction.

The block-size ablation, $B=8$ against $B=4$, is an honest trade and not a win. Pixel precision improves from 0.9599 to 0.9794, but in-region PSNR falls from 29.63~dB to 26.78~dB, because the payload drops from 128 to 32 bits per block. The tag halves from 32 to 16 bits, raising the per-block false-accept probability from $2^{-32}$ to $2^{-16}$, about 0.25 expected false accepts on a 16,384-block image. Finer localization is thus bought with weaker authentication and recovery.

\subsection{Limitations}
The scheme is fully fragile: any lossy re-encode, resize or filter destroys the watermark, so it protects images only inside a lossless custody chain. One-to-one allocation fails all-or-nothing per block, and reference sharing degrades more gracefully. Recovery is a low-frequency approximation from 12 DCT coefficients, not the original block. Pixel precision is capped by block quantization, the two highest corpus ratios are nominal only, and at $B=4$ the tag halves to 16 bits. Security rests on key secrecy. No claim of superiority over learned proactive schemes such as EditGuard \cite{editguard} is made.
